\subsection{Búsqueda binaria}

La búsqueda binaria intenta encontrar un valor específico, el \emph{objetivo}, en una secuencia ordenada. El problema que resuelve es entonces:
\begin{especificacion}
  \entrada{arreglo ordenado de enteros y un número entero: el objetivo}
  \salida{índice del objetivo en el arreglo o -1 si no está}
\end{especificacion}

Aplicando el paradigma, evidentemente el subproblema es buscar el objetivo solo en un segmento del arreglo. Imaginar que ese subproblema está resuelto equivale a saber en cuál segmento buscar (no a saber exactamente dónde está, porque entonces estaría resuelto el problema completo).

Implementar esa idea no es difícil: como el arreglo está ordenado, se puede comparar el objetivo con algún elemento del arreglo: si es mayor, necesariamente tendrá que estar a la derecha y entonces solo hace falta buscar en ese segmento; si es menor, tendrá que estar a la izquierda (y si es igual, ya no hace falta buscar más).

Lo más sensato es iniciar comparando con el elemento de la mitad, para que los dos posibles segmentos en los que tocaría buscar sean del mismo tamaño. Eso reduce el espacio de búsqueda a la mitad. Luego, se repite: una vez se sabe en qué segmento buscar, nuevamente se aplica la idea para reducir ese espacio a la mitad, hasta que se encuentra el elemento o se determina que no está.

\subsubsection{Solución recursiva}
\label{sec:busqueda_binaria_recursiva}

La implementación más natural es recursiva, porque permite representar muy naturalmente esa idea de buscar en solo una mitad del segmento cada vez. No hace falta modificar el arreglo para segmentarlo, basta con tener los índices del extremo izquierdo y derecho del segmento que se está tomando:

\begin{pseudocode}
def binary_search_rec(nums: arreglo[int], target: int, lo: int, hi: int) -> int:
    |\br{mi}|if hi < lo:
        return -1|\br{mf}|

    mid = |\bxl{over}|(hi - lo) // 2 + lo|\bxr{over}|

    |\br{gtb}|if target > nums[mid]:
        return binary_search_rec(nums, target, mid+1, hi)|\br{gtf}|
    |\br{ltb}|if target < nums[mid]:
        return binary_search_rec(nums, target, lo, mid-1)|\br{ltf}|

    |\bxl{eq}|return mid|\bxr{eq}|
\end{pseudocode}
\begin{annotations}
  \codebrace[width=11cm]{mi}{mf}{Si \inline{hi} está a la izquierda de \inline{lo}, los índices se cruzaron, lo que indica que el elemento no está.}
  \codebox[xshift=0.5cm, width=6cm]{over}{\annMitadSinOverflow{lo}{hi}}
  \codebrace[xshift=0.3cm, width=5cm, yshift=-0.1cm]{gtb}{gtf}{El objetivo es mayor: debe estar a la derecha. El llamado recursivo va desde después de la mitad hasta el final del segmento.}
  \codebox[xshift=0.5cm, width=12cm]{eq}{Si no es ni mayor ni menor, es igual: el objetivo está en el índice del medio.}
\end{annotations}

En los llamados recursivos se toma desde el índice de la mitad \emph{más uno} o \emph{menos uno}. Eso cumple dos funciones: primero, que no se vuelva a comparar con la mitad; pero sobre todo, permite que a media que el segmento del arreglo que se contempla es cada vez más pequeño, tal que el índice de la mitad se acerca más y más a algún extremo, eventualmente se crucen y la recursión llegue al caso base. Si solo se tomara el índice de la mitad, jamás se cruzaría con algún extremo y la recursión no terminaría cuando el objetivo no se encuentra.

Por último, la función recibe, además del arreglo y el objetivo, los índices. Así que técnicamente no recibe las entradas del problema propuesto. Envolverla en una función que sí las reciba es trivial:
\begin{pseudocode}
def binary_search(nums: arreglo[int], target: int) -> int:
    return binary_search_rec(nums, target, 0, len(nums) - 1)
\end{pseudocode}

\subsubsection{Complejidad}

El mejor caso para la búsqueda binaria es que el objetivo exista y esté en el centro del arreglo. En ese caso la ejecución es constante: \(T_\text{mejor} \in \Theta(1)\).

El peor caso es cuando se realizan el máximo número de pasos, que es cuando toca dividir el problema hasta que se que se llega a un problema de tamaño 1. La entrada de tamaño \(n\) se divide entre \(2\) una cantidad \(x\) de veces:

\begin{gather*}
  n \underbrace{{} \cdot \frac{1}{2} \cdot \frac{1}{2} \cdots \frac{1}{2}}_{x \ \text{veces}} = 1 \\
  \frac{n}{2^x} = 1
\end{gather*}

Despejando \(x\) para saber la cantidad de divisiones que se hace:

\begin{gather*}
  2^x = n \\
  x = \log_2(n)
\end{gather*}
Es decir, el número de veces que se divide el problema crece logarítmicamente con respecto al tamaño de la entrada. El resto de operaciones en el algoritmo son de tiempo constante, \(T_\text{peor} \in \Theta(\log n)\). 

Aunque ese peor caso ocurre cuando el objetivo no se encuentra, la complejidad para el caso promedio es la misma: \(T_\text{promedio} \in \Theta(\log n)\). Se omite la demostración matemática, pero la intuición es que los pasos de la búsqueda binaria se pueden diagramar como un árbol binario balanceado, y en todo árbol binario balanceado la mayoría de los nodos están cerca de las hojas (la parte de abajo es más gruesa). Por ende, en la mayoría de los casos, el algoritmo tendrá que ejecutar la máxima o cerca a la máxima cantidad de pasos.

Como \(T_\text{mejor} \in \Theta(1)\) y \(T_\text{promedio}, T_\text{peor} \in \Theta(\log n)\), búsqueda binaria es \(O(\log n)\).

\subsubsection{Solución iterativa}
\label{sec:busqueda_binaria_iterativa}

También se puede realizar una implementación iterativa de búsqueda binaria, en donde la forma de segmentar el arreglo consiste en desplazar el índice de algún extremo hacia la mitad:
\begin{pseudocode}
def binary_search_iter(nums: arreglo[int], target: int) -> int:
    lo = 0
    hi = len(nums) - 1

    while lo <= hi:
        mid = (hi - lo) // 2 + lo
        |\br{ib}|if target < nums[mid]:
            hi = mid - 1|\br{if}|
        elif target > nums[mid]:
            lo = mid + 1
        else:
            return mid

    return -1
\end{pseudocode}
\begin{annotations}
  \codebrace[xshift=2cm, width=8.5cm]{ib}{if}{El objetivo es menor que el valor del medio: debe estar a la izquierda. Se mueve el extremo derecho a justo antes de la mitad.}
\end{annotations}

El análisis asintótico de complejidad es el mismo de antes, pues se puede aplicar el mismo raciocinio para el máximo número de pasos. No obstante, esta implementación es ligeramente más eficiente al evitar el \hyperref[sec:sobrecosto_recursion]{sobrecosto de la recursión}.

\practicar[leetcode]{Binary Search}{https://leetcode.com/problems/binary-search/}
